% 350_Schwarze_Loecher_Dualitaet_De_ch.tex
% Johann Pascher, 4. September 2026

\providecommand{\BM}{\textbf{[B]}}
\providecommand{\KM}{\textbf{[K]}}
\providecommand{\SM}{\textbf{[S]}}

\phantomsection
\addcontentsline{toc}{chapter}{Schwarze Löcher und Zeit-Masse-Dualität}

\section*{Statuszeichen}
\begin{center}
\begin{tabular}{@{}lp{0.85\linewidth}@{}}
\textbf{[K]} & \emph{Konfirmiert} --- algebraisch vollständig. \\[4pt]
\textbf{[B]} & \emph{Belegt} --- durch Prüfskript gesichert. \\[4pt]
\textbf{[S]} & \emph{Spekulativ} --- Vermutung, noch offen. \\
\end{tabular}
\end{center}

\textbf{Hinweis.} Dieses Dokument ist als Ganzes \SM. Es zieht Konsequenzen
aus $T\cdot m=1$ (Dok.~003) und $G=\xi^2/(4m)$ (Dok.~012, R58) für ein
Regime --- schwarze Löcher, Hawking-Verdampfung --- das der Korpus bisher
nur am Rand berührt (Dok.~025: Hawking-Strahlung als Rückführung ins
$\xi$-Feld; Dok.~049: Feldknotentunneln durch den Horizont; Dok.~306:
Horizont als Extremum $T_0=\xi\,t_P$, $E_\text{max}=E_P/\xi$). Die
algebraischen Relationen sind durch \texttt{pruef\_350} gesichert; ihre
physikalische Gültigkeit in diesem Regime ist es nicht.

\section*{Abstract}

Aus $T\cdot m=1$ und $G=\xi^2/(4m)$ folgt $G\propto T$: die
Gravitationskopplung eines Objekts wächst mit seiner intrinsischen Zeit,
während $m\cdot G=\xi^2/4$ invariant bleibt \BM. Für ein verdampfendes
schwarzes Loch bedeutet das: die Sequenz der emittierten Muster ist durch
die Massenhierarchie festgelegt, und das letzte stabile Muster ist das
Neutrino --- der Zustand mit der längsten intrinsischen Zeit und
algebraisch das Vakuum des Clifford-Fock-Raums (Su[0]$=$Vss[0], Dok.~344)
\BM. Eine punktförmige Singularität mit endlichem $G$ ist ausgeschlossen
\BM\ (Dok.~078). Ob das Elektron-Muster (oder irgendein Muster) bei stark verändertem
$T$ noch dasselbe Objekt ist, bleibt offen \SM.

\section{Ausgangslage: die Frage}

Löst man die FFGFT-Gravitationsformel für eine Planck-Masse --- geht $G$
dann gegen Null bei endlicher Masse? Und gibt es eine berechenbare
Differenz zur Standardrechnung?

Die Antwort ist nicht ,,Ja'' oder ,,Nein'', sondern: die Frage setzt
voraus, dass $G$ und $m$ unabhängig variierbar sind. In FFGFT sind sie es
nicht.

\section{Satz A --- Kovarianz von $G$ unter $T\cdot m=1$}

\begin{theorem}[$G$ skaliert mit der intrinsischen Zeit \BM]
Aus $T\cdot m=1$ (Dok.~003) und $G=\xi^2/(4m)$ (Dok.~012) folgt
\[
  G = \frac{\xi^2}{4}\,T, \qquad m\cdot G = \frac{\xi^2}{4}=\text{const.}
\]
\end{theorem}

\begin{proof}
Einsetzen von $m=1/T$ in $G=\xi^2/(4m)$.
\end{proof}

\textbf{Lesart.} Wächst die intrinsische Zeit eines Objekts (z.\,B.\ durch
Zeitdilatation nahe einem Horizont), so sinkt seine Masse und steigt seine
Gravitationskopplung --- das Produkt beider ist eine Galois-Zahl. Ein
,,$G\to0$ bei endlicher Masse'' ist damit ausgeschlossen: $G\to0$ hieße
$T\to0$, also $m\to\infty$.

\section{Satz B --- Keine punktförmige Singularität}

\begin{theorem}[Massensingularität ausgeschlossen \BM]
Da $m\cdot G=\xi^2/4$ invariant ist, gibt es keinen Zustand mit
$m\to\infty$ bei endlichem $G$ und keinen mit $G\to\infty$ bei endlichem
$m$. Der klassische Singularitätsbegriff (endliche Masse in verschwindendem
Volumen bei konstantem $G$) hat in FFGFT kein Gegenstück.
\end{theorem}

\begin{proof}
Dok.~306 (R50) leitet den Ausschluss nativ her: aus $T\cdot m=1$ folgt
$E\to\infty\Leftrightarrow T\to0$; FFGFTs Zeit hat aber ein minimales,
nicht verschwindendes Inkrement $d_1=100\,\xi_1=1/75$, festgelegt durch
$\xi=4/30000$. Damit ist $T$ nach unten und $E$ nach oben beschränkt.
Der Horizont eines schwarzen Lochs ist das \emph{Extremum}: langsamste
Zeit $T=T_0=\xi\,t_P$, maximale Konzentration $E_\text{max}=E_P/\xi$, mit
exakter Schließung $T_0E_\text{max}=1$ --- alles endlich. Dok.~078 gibt
dieselbe Aussage als Ausschluss-Theorem ($m=\infty$ erzwingt $T=0$,
außerhalb des Definitionsbereichs). Mit $G=\xi^2/(4m)$ (Dok.~012) folgt,
dass $G\to0$ und $G\to\infty$ ebenso ausgeschlossen sind.
\end{proof}

\textbf{Hinweis zur Begründung.} Dok.~025/063 hatten den Ausschluss
ursprünglich über Heisenbergs Energie-Zeit-Unschärfe begründet; R50 hat
diese Ableitung durch die korpuseigene (Dok.~306) \emph{ersetzt} --- die
Unschärferelation ist hier nicht anwendbar (kein kanonischer
Zeitoperator, Pauli). Die Konklusion (kein singulärer Anfang, endlicher
Kern mit $\xi$-gesetzter Skala) ist unverändert. Konsistent mit R58
(keine separate Quantengravitation nötig, weil $T^4/Z_3$ pre-geordnet
ist). Was \SM\ bleibt, ist nur die Anwendung auf das Innere \emph{jenseits}
des Horizonts --- der Horizont selbst ist nach Dok.~306 das Extremum, kein
Übergang zu einer Singularität.

\section{Satz C --- Hierarchie der intrinsischen Zeiten}

\begin{theorem}[Neutrino als Zustand längster intrinsischer Zeit \BM]
Unter $T\cdot m=1$ gilt $T_\nu/T_e = m_e/m_\nu \approx 10^7$ und
$T_e/T_t = m_t/m_e \approx 3\cdot10^5$. Das Neutrino ist das massive
Fermion mit der längsten intrinsischen Zeit. Algebraisch ist es der
Vakuumzustand: Su[0]$=$Vss[0] (Dok.~344, Satz~C) \BM.
\end{theorem}

Dok.~025 gibt dieselbe Hierarchie als $\tau_n=\tau_\xi/\xi^n$: jede
Strukturebene (Teilchen, Atom, Molekül, Zelle) hat eine um $1/\xi$ längere
Zeitskala. Das Neutrino sitzt am unteren Ende der Massenskala und damit am
oberen Ende der Zeitskala.

\section{Satz D --- Verdampfungssequenz}

\begin{theorem}[Ordnung der Hawking-Emission nach Masse \SM]
Ein verdampfendes schwarzes Loch verliert Masse; mit $M\to0$ wächst
$T_\text{Loch}$. Die emittierbaren Muster sind bei jedem $T$ diejenigen,
deren eigene intrinsische Zeit zu $T_\text{Loch}$ passt. Die Sequenz folgt
der Massenhierarchie:
\[
  t\to b\to\tau\to c\to\mu\to s\to d\to u\to e\to\nu.
\]
Das letzte stabile Muster vor dem vollständigen Auflösen ist das
Neutrino.
\end{theorem}

\textbf{Verbindung zum Korpus.} Dok.~025 nennt Hawking-Strahlung als
,,Energie-Rückführung ins $\xi$-Feld''. Satz~D präzisiert: die Rückführung
läuft durch die Massenhierarchie hindurch und endet beim Zustand, der
algebraisch bereits das Vakuum ist. Dok.~049 beschreibt denselben Vorgang
als Feldknotentunneln durch den Horizont --- Satz~D ordnet die Knoten nach
ihrer intrinsischen Zeit.

\section{Was mit einem Elektron nahe dem Horizont geschieht \SM}

Nahe dem Horizont ist die lokale Zeit gedehnt. Unter $T\cdot m=1$ heißt
das: $m_e$ sinkt, $G_e$ steigt. Aber $m_e$ ist nicht isoliert ---
$\alpha=\xi E_0^2/K_\text{frak}$ (Dok.~011), $v=m_e/((4/3)\xi^{3/2})$
(Dok.~006), $G_F$, alle Yukawa-Kopplungen hängen an derselben Skala. Ändert
sich $T$, ändern sich alle kohärent.

Damit ist die Frage ,,was passiert mit dem Elektron'' nicht beantwortbar,
ohne zu klären, ob das Elektron-\emph{Muster} --- Galois-Orbit $O_1$,
QR-Klasse, Sd[3]-Zustand (Dok.~346--349) --- bei verändertem $T$ noch
dasselbe Objekt ist. Der Orbit ist eine algebraische Invariante; die
Masse, die ihm zugeordnet ist, ist es nicht. Ob man das Ergebnis noch
,,Elektron'' nennen soll, ist eine Definitionsfrage, die der Korpus nicht
entscheidet.

Was als Hawking-Strahlung austritt, ist thermisch --- ein Mischzustand, in
dem die Identität des eingefallenen Musters nicht mehr als benanntes
Teilchen, sondern als Beitrag zum Spektrum vorliegt. Das
Informationsparadoxon wird dadurch nicht gelöst, sondern in die
$T$-Struktur des Endzustands verschoben \SM.

\section{Antwort auf die Ausgangsfrage}

\begin{enumerate}
\item $G\to0$ bei endlicher Masse: \textbf{nein} --- $G\propto T$, und
  $T\to0$ erzwingt $m\to\infty$ \BM.
\item Berechenbare Differenz zur Standardrechnung: \textbf{ja}, aber nicht
  als Zahlenwert von $G$ (der ist lokal identisch, Dok.~012), sondern als
  Struktur: $m\cdot G=\xi^2/4$ ist eine Galois-Zahl, und die Planck-Skala
  ist die Skala, auf der die \emph{lokalen} Werte von $G$, $\hbar$, $c$
  zusammenpassen --- kein universeller Grenzwert \SM.
\item Was am Ende bleibt: ein neutrino-artiges Muster, dann das Vakuum
  Vss[0] \SM.
\end{enumerate}

\section{Zusammenfassung}

\begin{table}[H]\centering
\caption{Ergebnisse Dok.~350}
\begin{tabular}{lll}\toprule
Satz & Inhalt & Status \\\midrule
A & $G=\tfrac{\xi^2}{4}T$; $m\cdot G$ invariant & \BM \\
B & Keine Massensingularität; Horizont $=$ Extremum (Dok.~306/078) & \BM \\
C & Neutrino $=$ längste intrinsische Zeit $=$ Vakuum & \BM \\
D & Verdampfungssequenz endet beim Neutrino & \SM \\
--- & Identität des Elektrons bei verändertem $T$ & offen \SM \\\bottomrule
\end{tabular}
\end{table}

\section*{Prüfskript}
\texttt{pruef\_350\_schwarzes\_loch.py} --- Sätze A--D, algebraische
Konsistenz bestanden. Physikalischer Gesamtstatus: \SM.

\section*{Vermerk zur Verwendung von KI-Assistenz}
Erstellt mit Unterstützung von Claude (Anthropic).
Physikalische Interpretationen und Schlussfolgerungen:
Johann Pascher (ORCID 0009-0000-6518-4064).

\begin{thebibliography}{9}
\bibitem{P003} J.~Pascher, \textit{Dok.~003: T0-Grundlagen}, FFGFT-Korpus.
\bibitem{P012} J.~Pascher, \textit{Dok.~012: Gravitationskonstante}, FFGFT-Korpus.
\bibitem{P025} J.~Pascher, \textit{Dok.~025: T0-Kosmologie}, FFGFT-Korpus.
\bibitem{P078} J.~Pascher, \textit{Dok.~078: Zeit --- Ausschluss-Theorem $T\cdot m=1$}, FFGFT-Korpus.
\bibitem{P306} J.~Pascher, \textit{Dok.~306: Native Zeit-Energie-Reziprozität}, FFGFT-Korpus (R50).
\bibitem{P049} J.~Pascher, \textit{Dok.~049: Lagrangian-Vergleich}, FFGFT-Korpus.
\bibitem{P344} J.~Pascher, \textit{Dok.~344: Furey-Fermionen in GALG und FFGFT}, FFGFT-Korpus (Sept.\ 2026).
\bibitem{P349} J.~Pascher, \textit{Dok.~349: SU$(2)_L$-Chiralität aus Galois}, FFGFT-Korpus (Sept.\ 2026).
\bibitem{Hawking1975} S.~W.~Hawking, Commun.\ Math.\ Phys.\ 43 (1975) 199.
\end{thebibliography}
